\documentclass[11pt]{article}
\usepackage[utf8]{inputenc}
\usepackage[T1]{fontenc}
\usepackage{lmodern}
\usepackage[margin=1in]{geometry}
\usepackage{csquotes}
\usepackage[style=authoryear,backend=biber,sorting=nyt,maxbibnames=99,giveninits=true,uniquename=init]{biblatex}
\usepackage[hidelinks]{hyperref}
\addbibresource{references.bib}
\begin{document}
\section*{Unintended Consequences and Multi-Objective Tradeoffs}

Subsurface drainage and the reductive treatment systems appended to it are
increasingly recognised for effects that fall outside the nitrate- and
yield-centred accounting that has historically justified them. Removing a
nutrient from one pathway can shift it to another, treating drainage water can
manufacture new by-products, and altering the field water regime can set slow
biogeochemical and pedological changes in motion. This section consolidates
these cross-constituent consequences---greenhouse-gas coupling, pollution
swapping, emerging contaminants, trace-element legacies, and drainage-induced
soil evolution---as a single frontier defined less by settled magnitudes than
by unmeasured tradeoffs. The recurring message is that few studies have
quantified more than one consequence at a time, so the net balance among
nitrate reduction, greenhouse-gas emission, secondary pollutant release, and
long-term soil change remains largely conjectural.

\subsection{Greenhouse-gas coupling and carbon export}

Greenhouse-gas (GHG) coupling is a conspicuous gap in drainage accounting.
Fluxes of carbon sequestration, pre-implementation emissions, and dissolved
nitrous oxide in tile drainage and conservation-practice outflows need
quantification \parencite{Mitchell2023}, and coordinated large-scale studies
are required because few have simultaneously measured N\textsubscript{2}O,
CH\textsubscript{4}, and soil-organic-carbon change at the same sites
\parencite{Hall2025}. Longer-term dynamics of wetland carbon storage,
N\textsubscript{2}O, and CH\textsubscript{4}, and the GHG effects of upgrading
drainage infrastructure, remain unresolved \parencite{Mitchell2022}, and carbon
budgets should be integrated into wetland-design models \parencite{Kadlec2012}.

A central strand couples nitrate leaching to indirect N\textsubscript{2}O.
\textcite{Sawamoto2003} call for study of the nitrate--N\textsubscript{2}O
relationship in groundwater and of indirect N\textsubscript{2}O from subsurface
drainage, \textcite{Minamikawa2015} argue for new calculation methods for
indirect N\textsubscript{2}O from leaching and runoff, and the tradeoff between
indirect N\textsubscript{2}O and heterotrophic-denitrification N removal must be
evaluated across groundwater, drainage, river, and estuary
\parencite{Minamikawa2013,Minamikawa2010}. Open questions include
nitrate--N\textsubscript{2}O linkages when N application and precipitation are
temporally disconnected \parencite{OBrien2022}, the magnitude and fate of
N\textsubscript{2}O losses \parencite{Burton2024}, whether dissolved
N\textsubscript{2}O losses exceed surface emissions \parencite{Smith2019}, and
long-term cumulative dissolved-GHG effects of management \parencite{Oliveira2025}.
\textcite{Webb2021} outline an emissions-accounting agenda---seasonal
emission-factor impacts, process-based modelling, gas-transfer-velocity models,
the DOC/NO\textsubscript{3} ratio, and a global database---while
\textcite{Fontes2019} recommend year-round monitoring of both N\textsubscript{2}O
and nitrate leaching. Net-GHG evaluation in longer experiments
\parencite{PrezaFontes2023}, the fate of waterborne methane alongside phosphorus
balances when selecting measures \parencite{Carstensen2021}, comparative
life-cycle assessment from complete N balances \parencite{HernandezRamirez2011},
and biogeochemical measurement and modelling of methane production and oxidation
in cold, poorly drained soils \parencite{Kumar2014b} extend the frontier.

Dissolved-organic-carbon (DOC) dynamics and their ecological implications remain
largely unknown \parencite{Ruark2009,Furukawa2008}, with dissolved-organic-matter
quality, microbial availability, and episodic drainage influence demanding
paired-watershed designs and scale-appropriate sampling
\parencite{Dalzell2011,Bellmore2015}. Treatment structures add their own GHG and
carbon budgets: \textcite{Brunton2024} recommend fuller assessment of full-size
bioreactor methane dynamics, age effects on GHG consumption and release,
``hot moment'' production, and real-time control on outflow redox potential, and
\textcite{Strock2025} note that a 2-h hydraulic residence time raises GHG
emissions and that DOC export needs monitoring. Carbon-export accounting more
broadly is a monitoring gap---aqueous total-carbon export for accurate
sequestration accounting, multi-year data under variable climate, and the
potential of water-level control to minimise dissolved-inorganic-carbon loss
\parencite{Schilling2023}.

\subsection{Pollution swapping and by-product formation in treatment systems}

Reductive treatment introduces tradeoffs that are only partly quantified.
\textcite{Christianson2021} name reducing pollution swapping---especially at
start-up---and improving understanding of the removal of additional pollutants,
compounds, and biologicals as key research areas. The by-products of concern
span the reduction sequence. Nitrite accumulation should be measured in future
bioreactor studies alongside long-term evaluation of phosphorus transport
\parencite{Bell2015}. Mercury methylation is a monitoring priority that
\textcite{Maxwell2022} suggest minimising through an outflow
nitrate-concentration design criterion rather than a mass metric. Redox-driven
phosphorus release is a recurring gap: \textcite{Bailon2025} call for evaluating
start-up flushes with immediate monitoring, aging bioreactors, and total
phosphorus alongside higher-frequency dissolved reactive phosphorus (DRP)
monitoring and drawdown plates, and \textcite{Rivas2020} for continued
multi-season monitoring to confirm longer-term DRP removal and the factors
governing it. The mechanisms of enhanced phosphorus uptake by woodchips under
anoxic--oxic cycling---including polyphosphate-accumulating organisms and whether
anaerobic conditions are necessary---remain to be elucidated
\parencite{Zhang2024b}. \textcite{Strock2025} tie these together, calling for
work on P-removal-material longevity, cold-condition optimisation, DOC-export
monitoring, and strategies for managing P-treatment materials to minimise
unintended phosphorus release, while \textcite{Schilling2023} note that
water-level control may curb dissolved-inorganic-carbon loss. Across these
by-products the common prescription is the same: longer, higher-frequency,
multi-constituent monitoring to bound tradeoffs that single-season demonstrations
cannot resolve.

\subsection{Emerging contaminants, pesticides, pathogens, and antibiotic resistance}

Emerging contaminants and biological interactions broaden the constituent scope
of drainage assessment. Priorities include better constraint of the mechanism
and lifetime of glyphosate in soil and water \parencite{Spiese2023}, planter
modifications to reduce aquatic exposure to neonicotinoid and other seed-applied
pesticides \parencite{Schaafsma2019}, how watershed microbial communities respond
to and control nutrient fluxes \parencite{Knorr2023}, and compound-source
identification across rainfall-runoff events \parencite{Fraindová2022}.
Particulate silicon transport coupled to warming and diatom dynamics
\parencite{Sethna2022}, nutrient stoichiometry and organic-matter fate in aquatic
ecosystems \parencite{Manninen2023,Manninen2018}, and release of accumulated
compounds and resistance genes in sludge-amended soils \parencite{Ghirardini2019}
extend the list.

\subsubsection{Contaminant transport through inlets and macropores}

How contaminants enter and move through drainage networks is explicitly
under-studied. \textcite{Li2017b} recommend research on how nitrate, phosphorus,
and pesticides move through drainage inlets, and on inlet performance under
crop-specific debris. The effective contributing area may differ for water,
tracers, and reactive contaminants \parencite{Akay2008}, and controls on
tile-sediment flocculation remain unresolved \parencite{Stone1997}. Pesticide
transport modelling forms a coherent cluster: particle-mediated colloidal
algorithms and separate macropore accounting for leaching \parencite{Wendell2024},
transport by mobile particles in convection--dispersion models \parencite{Zehe2001},
a quasi-two-dimensional profile distribution and macropore--drain interconnectivity
in chemical routing \parencite{Fox2007b}, slow-immobilisation and
dissociation-kinetics processes for persistent pesticides \parencite{Novak2003},
multi-season field monitoring across soil types \parencite{Dousset2004}, and
vulnerability functions from field characteristics \parencite{Leu2004}. Analytical
and marker methods---hydroxy-triazine determination \parencite{MUIR1978}, a dicamba
marker \parencite{Heaven2012}, continued linuron monitoring \parencite{Gordon2005},
and over-winter freezing/thawing effects on herbicide fate \parencite{Siimes2006}
---complete this cluster.

\subsubsection{Pathogens, microbial ecology, and antibiotic resistance}

Microbial contaminant transport spans several open questions. Pathogen removal
under lower residence times, \textit{Salmonella} and \textit{E.\ coli} removal
mechanisms, and behaviour at higher flows need investigation
\parencite{Soupir2018}. The ecology of indicator bacteria in treatment
wetlands---survival of faecal versus cryptic ``\textit{E.\ coli}-like'' strains
\parencite{Stott2023} and curli expression and downstream risk
\parencite{Truhlar2015}---is unresolved, as is long-term pathogen leaching after
manure application \parencite{Malik2004}. Fecal-source tracking needs
corroborating markers \parencite{Gottschall2013}, irrigation with wastewater
warrants disinfection and reduced cold-season rates \parencite{Ding2021}, and a
broad catchment-scale microbial agenda spans riparian entrapment, constructed
wetlands, and water as a human-infection vector \parencite{Collins2007}.
Antibiotic-resistance genes are an explicitly emerging frontier: their diversity,
persistence, and dissemination require qPCR to complement metagenomics
\parencite{Smith2019b}, and further study is needed to link resistance genes to
host identity and transfer mechanisms \parencite{Rieke2018}.

\subsection{Trace-element legacies and ecotoxicology}

Trace-element contamination associated with drainage in the western United States
represents a long-horizon monitoring commitment. \textcite{BENSON1993} outline
continued biological and groundwater monitoring, research into 10--20-year
selenium-dissipation strategies, annual management planning, and pilot-scale
microbial-volatilisation demonstration, and \textcite{SAIKI1987} note further
studies to determine whether aquatic animals elsewhere in the San Joaquin Valley
are contaminated by selenium. Ecotoxicological monitoring of other contaminants
extends the frontier: \textcite{Campbell2023} call for mechanistic work on
reduced amphibian cerebellum size, field examination of imidacloprid and its
olefin metabolite, and regular monitoring of neonicotinoids near agricultural
waterways with tighter restrictions to protect non-target organisms, while
\textcite{Dreyfuss2018} note that effects on \textit{Galba truncatula} raise
conservation questions given its role as intermediate host of \textit{Fasciola
hepatica}. These constituents share a common feature: their consequences unfold
over decades and reach organisms far removed from the field, so short monitoring
windows systematically understate them.

\subsection{Drainage-induced soil evolution, health, and biology}

Slower-acting soil and biological responses to drainage are repeatedly flagged as
understudied. The soil-biology response to changing drainage hydrology requires
longer-term evaluation of biological indicators under standardised sampling
\parencite{Sherbine2023}, study of interactions between subsurface drainage and
soil physicochemical properties and microbial activity \parencite{Yang2025}, and
continued research on microbial community structure and function to explain
tillage and drainage impacts \parencite{Kaur2024b,Kumar2014}. Robust soil-health
diagnostics are sought---scalable nematode measurement within a soil-health index
\parencite{Li2023b}, valid and repeatable soil-quality indicators in a systems
framework \parencite{Nakajima2015}, soil-enzyme activities as management
diagnostics \parencite{Singh2007}, proof that no-till reduces nitrate leaching
sustainably \parencite{Masri2024}, and dynamic compaction indicators with
long-term restoration trials \parencite{Caron2025}. Specific biological questions
include \textit{Fusarium} abundance and soybean root rot under long-term drainage
\parencite{Han2017}, disruption of crack-to-burrow linkages by conservation
tillage \parencite{Shipitalo2004}, evapotranspiration quantification to explain
drain-flow differences \parencite{Daigh2014}, cultivar drought tolerance guiding
system selection \parencite{DarziNaftchali2017b}, and forage-sorghum nitrate
testing \parencite{Nash2017}.

A distinctive line treats subsurface drainage as a driver of long-term,
potentially chaotic soil evolution. \textcite{Montagne2009} argue that
quantifying human-induced soil evolution as a function of time---including
drainage's impact on clay-mineral dissolution---should be a research priority,
and that soil evolution under climate change must be explicitly studied and
incorporated in global models \parencite{Montagne2010}; significant material
losses over a few decades make anthropogenic perturbation a powerful tool for
studying soil-evolution timescales \parencite{Montagne2013}. Whether
drainage-induced evolution continues or stabilises over decades, resolved through
measurements along drains at different slope positions and network ages, remains
open \parencite{Frison2009}, and a long classical experiment continues to test
whether soil structure is eventually affected \parencite{HOOGHOUDT1952}.
Empirical prediction of earthworm activity from tillage, residue cover, and
particle size is proposed \parencite{Ouellet2008}, and sustained drainage
operation and cultivation is framed as a management imperative for long-term soil
health \parencite{Singh2007}. Collectively, these consequences reframe drainage
as an agent of slow, coupled biogeochemical and pedological change whose full
cost cannot be read from a few years of outlet monitoring.

\printbibliography
\end{document}
